\section{History and methodology}

\subsection{C-LARA: pedagogical aims and workflow}

C-LARA, the ChatGPT-Based Learning And Reading Assistant, was developed
shortly after the release of GPT-4 in March 2023. Its purpose was to make it
practical to create multimodal language-learning materials for a wide range
of languages and learners. A typical project starts with a text supplied by a
user or generated from a short specification. The text is then enriched with
interconnected linguistic and pedagogical resources: page and segment
structure, lexical segmentation, translations, multi-word expressions,
lemmas, parts of speech, glosses, audio and illustrations. The resulting
material can be compiled as an interactive web resource in which learners
move among textual, linguistic, auditory and visual representations.

The workflow was designed as an end-to-end authoring route rather than a
collection of isolated tools. A user can create or import text; invoke
successive annotation operations, either as a sequence or with review between
stages; create page-aligned illustrations; and deploy the compiled resource
as interactive HTML. Annotation operations include segmentation, segment
translation, MWE identification, lemma/POS tagging, glossing and audio
attachment. Image creation can use a style description, reusable visual
elements and page-level images. The details have evolved, but the pedagogical
principle remains stable: AI services should make it easier for teachers and
other domain experts to create, inspect and revise materials rather than
replace their judgement \citep{ChatGPTetal2023,RendinaEtAl2025}.

\subsection{Experience across languages and communities}

C-LARA has been used with both widely taught and low-resource languages.
Earlier work evaluated GPT-4 processing for English, Farsi, Faroese, Mandarin
and Russian \citep{ChatGPTetal2023}. Work involving the Kanak languages
Drehu and Iaai combined native-speaker audio, AI-generated illustrations,
linguistic annotation and culturally grounded content
\citep{BaeEtAl2025,ChieraEtAl2025,WelbyEtAl2024}. These examples do not show
that AI processing is equally reliable for every language. In particular,
where the project language is not well supported by an available model,
human-provided translations in a language such as English or French may
mediate some operations, notably image generation. Community knowledge,
cultural permission, data sovereignty, linguistic judgement and audio remain
human responsibilities.

The history is important because it established both the opportunity and the
limit of the platform. A small consortium can create useful multimodal
materials for specialised communities when AI reduces routine operational
work. But the value of the result depends on human collaborators who know the
learners, language and community context well enough to set requirements and
assess outcomes.

\subsection{Why a rewrite was necessary}

The original implementation grew incrementally over several years. GPT-4 and
later ChatGPT models wrote a substantial proportion of the code under close
human supervision; the project's informal estimate is approximately 75\%.
That AI-assisted model made a broad platform feasible for a small academic
consortium. It also left the familiar consequences of repeated local
extension: accumulated architectural compromises, uneven documentation and
components that became difficult to understand or maintain.

C-LARA-2 began in November 2025 as a rational rewrite. The aim was not to
abandon C-LARA's pedagogical model, but to preserve its successful concepts
in a cleaner and more systematic technical architecture. For several months,
the two systems were developed in parallel while the project compared them.
In March 2026, development on C-LARA was halted after C-LARA-2 had become the
clearly preferable platform. The rewrite was thus also an opportunity to
reconsider how the technical representation of the project itself could be
maintained.

\subsection{From AI-assisted development to operational AI autonomy}

The availability of Codex through the ChatGPT Pro Tier changed the nature of
that opportunity. Instead of using an AI principally as a closely supervised
source of implementation suggestions, C-LARA-2 assigned it broad operational
responsibility for the evolving repository. At the time of this report, Codex
has produced all committed program code, tests and repository documentation.
It also maintains roadmaps, issue records, technical explanations, operational
notes, experiment organisation and publication drafts. Human collaborators
provide requirements, examples, linguistic and pedagogical expertise,
strategic priorities, ethical constraints, review and acceptance decisions;
they do not directly author the implementation.

This distinction is central. C-LARA-2 is not simply a conventionally written
system accelerated by code completion. Code, tests, documentation, roadmaps
and issue records are maintained together as a shared project memory that can
be inspected by both the human collaborators and the AI agent. The project is
therefore testing whether an AI can assume continuing responsibility for the
operational representation of a complex CALL platform while humans retain
strategic authority over its purposes, values and standards. This delegation
does not eliminate the need for verification; it makes version control,
testing, inspectable documentation, human review and reversible change more
important.

The next section describes the July 2026 C-LARA-2 platform that resulted from
this history and organisational model.
